\section*{Chapitre \ref{ch:nw:fx-connectivity} --- Connectivité FX}

\begin{solution}{exo:nw:fx-connectivity:1}
Le FX spot de New York (et les métaux) se trouve chez Equinix Secaucus NY5, avec une sauvegarde à Slough LD4.2~; le FX spot et les NDF de Londres sont en LD4.2, avec une sauvegarde
dans NY5. Chaque produit se négocie en un seul lieu.
\end{solution}

\begin{solution}{exo:nw:fx-connectivity:2}
Seuls les niveaux de prix où la firme dispose d'un crédit bilatéral auprès des contreparties~; la quantité à un niveau peut inclure des ordres
qu'elle ne peut pas exécuter, si bien qu'elle peut dépasser ce que la firme peut exécuter.
\end{solution}

\begin{solution}{exo:nw:fx-connectivity:3}
\qty{46.85}{\milli\second} dans un sens~; à 1.5, \qty{70.3}{\milli\second} dans un sens et \qty{140.5}{\milli\second} en aller-retour.
\end{solution}

\begin{solution}{exo:nw:fx-connectivity:4}
Le fournisseur est à l'origine~: sa cotation part vers chaque site dès qu'il a valorisé le mouvement, sur un chemin aussi rapide que celui du preneur de liquidité, de sorte que
la seule différence tient aux \qty{0.1}{\milli\second} de valorisation.
\end{solution}

\begin{solution}{exo:nw:fx-connectivity:5}
\qty{79.5}{\milli\second}~: le réseau à 1.5 plus \qty{40}{\micro\second} d'étapes. Un serveur à Secaucus, à côté du moteur, économiserait
presque tout ce délai, mais apprendrait les nouvelles de Tokyo avec \qty{79.5}{\milli\second} de retard, de sorte qu'il n'aide que pour les négociations dirigées par New York.
\end{solution}

\begin{solution}{exo:nw:fx-connectivity:6}
Après un mouvement, les cotations qui n'ont pas encore suivi sont les meilleures du côté que le marché a quitté, et elles proviennent des fournisseurs
les plus éloignés du mouvement~; l'agrégateur les sélectionne parce qu'elles paraissent les meilleures, précisément pendant qu'elles sont périmées.
\end{solution}

\begin{solution}{exo:nw:fx-connectivity:7}
\qty{47}{\milli\second}~: 99.4~\% sont détectées~; à \qty{46}{\milli\second}, 96.2~\%.
\end{solution}

\begin{solution}{exo:nw:fx-connectivity:8}
Ce n'est pas le caractère périmé que le contrôle doit couvrir~: avec des chemins égaux, le contrôle de prix à l'arrivée connaît déjà le mouvement. Si les exécutions sont toxiques,
le chemin d'information du fournisseur est plus lent que ceux des preneurs de liquidité~; une fenêtre plus longue que ne l'exigent les contrôles va à l'encontre du Principe 17. La correction
consiste à accélérer le chemin d'information ou à valoriser à Tokyo.
\end{solution}

\begin{solution}{pb:nw:fx-connectivity:1}
\textbf{Partie I.}
\begin{enumerate}
\item \qty{27.07}{\milli\second} et \qty{46.85}{\milli\second}.
\item \qty{40.6}{\milli\second} et \qty{70.3}{\milli\second}.
\item \qty{140.6}{\milli\second} à TY3, \qty{31.5}{\milli\second} à NY5 et NY6, \qty{111.1}{\milli\second} à SG1, et
  \qty{0.1}{\milli\second} à LD4, où le preneur de liquidité n'en est pas informé plus tôt que le fournisseur.
\item \qty{81.3}{\milli\second} à NY5 et NY6, \qty{31.5}{\milli\second} à TY3, \qty{8.2}{\milli\second} à SG1 et
  \qty{0.1}{\milli\second} à LD4.
\end{enumerate}
\textbf{Partie II.}
\begin{enumerate}[resume]
\item Aucun délai supplémentaire au-delà du temps de valorisation~: \qty{0.1}{\milli\second}.
\item 70~\% immédiatement et 99.4~\% après \qty{2}{\milli\second}.
\item \qty{185.6}{\milli\second} et \qty{45.1}{\milli\second}.
\item \qty{47}{\milli\second}.
\end{enumerate}
\textbf{Partie III.}
\begin{enumerate}[resume]
\item Un contrôle de validité, incluant un crédit disponible suffisant, et un contrôle de prix par rapport au prix courant disponible pour le client.
\item Non~: un délai supérieur à ce qu'exigent les contrôles va à l'encontre du Principe 17.
\item Un chemin d'information plus rapide depuis Tokyo, ou valoriser les mouvements de Tokyo depuis un serveur à Tokyo~; des contrôles communiqués et prompts.
\item À TY3, où débutent les flux vers le site de Tokyo~: les cotations de ce site seraient fraîches après les mouvements de Tokyo, et périmées après ceux
  de Londres et de New York à la place.
\end{enumerate}
\textbf{Partie IV.}
\begin{enumerate}[resume]
\item \emph{Résultat nommé}~: la cotation USDJPY d'un fournisseur de Londres est périmée pendant \qty{140.6}{\milli\second} sur le site de Tokyo après un mouvement de Tokyo, avec
  tous les chemins à 1.5 fois le plancher, et \qty{185.6}{\milli\second} s'il apprend les mouvements de Tokyo par un backbone à 2.46~; le temps de rétention qui le neutraliserait
  ne serait que le temps de valorisation avec des chemins égaux, et \qty{45.1}{\milli\second} (47 pour détecter 99~\% avec la gigue (jitter)) avec le backbone, un
  délai que le Code n'autorise pas, si bien que le remède porte sur le chemin et non sur la fenêtre.
\item Publié~: les sites, le texte du Code, les \omterm{def:nw:the-north-american-map:factor}{facteurs de tracé} du chapitre~19. Supposé~: les facteurs appliqués à ces routes, le temps de valorisation, la
  gigue (jitter), et l'endroit où naît la nouvelle.
\item La cotation du fournisseur de Tokyo est fraîche, celle du fournisseur de Londres est périmée pendant \qty{140.5}{\milli\second}, celle du fournisseur du New Jersey pendant
  \qty{158.9}{\milli\second}~; il sélectionne les cotations périmées du côté que le marché a quitté.
\item Horodater chaque cotation à sa création et à chaque accusé de réception ou écho des données de marché de la plateforme, et comparer les exécutions avec son propre mid
  un court instant plus tard, par site et par origine du mouvement.
\item Ses ordres sont transmis à New York ou à Londres, de sorte qu'il négocie au temps du moteur, à un saut de réseau près, quelle que soit sa passerelle.
\item Ils raccourcissent les chemins des preneurs de liquidité~: le chemin d'information du fournisseur est alors plus lent que les leurs, et le temps de rétention nécessaire augmente de la
  différence, là encore une affaire de chemins et non de fenêtres.
\item Dans les lignes FX~: les plateformes et les sites par paire, les sites et les chemins de valorisation du fournisseur, le crédit, et le suivi des cotations périmées.
\item Si le fournisseur apprend un mouvement aussi vite que la demande du preneur de liquidité peut l'atteindre.
\end{enumerate}
\end{solution}

\begin{solution}{iq:nw:fx-connectivity:1}
À LD4 à Slough et sur le campus de Secaucus, où se trouvent les principaux moteurs et flux, avec une présence à TY3 pour les clients de Tokyo et
streaming~; SG1 pour les NDF. Chaque paire s'apparie en un seul lieu, de sorte que les serveurs se trouvent là où les paires que vous négociez s'apparient et où naissent leurs nouvelles.
\iqlookfor{Moteurs par paire~; les campus de Slough et de Secaucus~; Tokyo comme accès et streaming.}
\end{solution}

\begin{solution}{iq:nw:fx-connectivity:2}
Un contrôle de risque portant sur la validité, y compris le crédit, et sur le prix, appliqué sans délai~; ce n'est pas une option gratuite, ni une source d'information pour le trading
contre le client, ni un moyen d'ajouter un délai au-delà des contrôles.
\iqlookfor{Principe 17~; validité et prix~; pas de délai ajouté~; transparence.}
\end{solution}

\begin{solution}{iq:nw:fx-connectivity:3}
Horodater les cotations à leur création sur une horloge synchronisée, lire les données de marché ou les accusés de réception des plateformes avec leurs horodatages sur chaque site,
et comparer avec l'instant où le mouvement qui aurait dû modifier la cotation a atteint chaque site par le chemin le plus rapide.
\iqlookfor{Synchronisation des horloges~; mesure par site~; le chemin le plus rapide comme référence.}
\end{solution}

\begin{solution}{iq:nw:fx-connectivity:4}
Si le chemin d'information du fournisseur a changé (une route, un prestataire, un commutateur), si les preneurs de liquidité ont gagné un chemin plus rapide, les rejets par motif,
et les exécutions comparées au mid quelques millisecondes plus tard, par site.
\iqlookfor{Chemins des deux côtés~; markouts~; motifs de rejet.}
\end{solution}

\begin{solution}{iq:nw:fx-connectivity:5}
C'est le maximum parmi les cotations périmées et fraîches, et après un mouvement les périmées sont les meilleures~; avec le crédit bilatéral et le last look, le meilleur
prix affiché peut aussi être un prix que le client ne peut pas négocier ou qui sera rejeté.
\iqlookfor{Sélection des cotations périmées~; filtrage du crédit~; rejets.}
\end{solution}

\begin{solution}{iq:nw:fx-connectivity:6}
Moteurs de valorisation à Londres, Secaucus et Tokyo, alimentés chacun par les chemins les plus rapides en provenance des autres sites et des principaux marchés des paires~; flux depuis
le moteur le plus proche de chaque site~; un service de crédit unique et cohérent entre les plateformes~; des contrôles de prix et de validité prompts et communiqués~; la péremption et
les markouts mesurés par site.
\iqlookfor{Valorisation au plus près de la nouvelle~; parité des chemins~; crédit~; suivi.}
\end{solution}
